% =====================================================================
% LES DEUX NOTES DE SYNTHÈSE, FRANÇAISE PUIS ANGLAISE
%
% PIÈCES PROPRES À LA VERSION ENSAE. Seul main_ensae.tex les appelle. Les
% consignes de l'école les exigent, autonomes, d'une à deux pages chacune, et
% lisibles par un non-spécialiste ; elles imposent aussi leur PLACE dans le
% fichier unique déposé, en toute fin de document. L'Institut des Actuaires ne
% demande ni l'une ni l'autre : le mémoire porte un résumé et un abstract, qui
% sont d'un autre genre, plus techniques et plus courts.
%
% Texte repris de l'ancien rapport de stage raccourci
% (exploratory/rapport_ensae/rapport_ensae.tex), où il avait été validé. Aucun
% nombre n'y a été modifié.
% =====================================================================

\chapter*{Note de synthèse}
\addcontentsline{toc}{chapter}{Note de synthèse (français)}
% UN \chapter* NE MET PAS LA MARQUE DE TITRE COURANT A JOUR, et fancyhdr garde
% alors celle du dernier \chapter numerote : les deux notes s'imprimaient sous
% un titre courant << Bibliographie >>. Vu sur le PDF rendu, pas sur la source.
\markboth{Note de synthèse}{}

% HARNAIS-HORS-SECTION: note de synthese autonome, destinee a un lecteur non
% specialiste. Elle reprend, en les arrondissant, des grandeurs etablies et
% verifiees dans le corps du memoire, ou elles portent leurs citations de script.

\textbf{Le problème.} Depuis janvier 2025, le règlement européen DORA impose aux
banques et aux assureurs la maîtrise de cinq domaines de résilience informatique :
gouvernance du risque technologique, gestion des incidents, tests de résilience,
maîtrise des prestataires tiers, partage d'informations sur les menaces. Le
règlement dit ce qu'il faut mettre en place. Il ne dit pas ce que coûte de ne pas
le faire, et le cadre prudentiel qui fixe le capital d'un assureur,
Solvabilité~II, ne le dit pas non plus : sa charge de risque opérationnel est un
pourcentage des primes, identique pour un assureur exemplaire et pour un assureur
défaillant. Un dirigeant qui arbitre un budget de conformité ne dispose donc
d'aucun chiffre.

\textbf{La difficulté, qui n'est pas celle qu'on croit.} Les cinq domaines ne sont
pas des cases indépendantes. Une gouvernance défaillante retarde la détection des
incidents, et un prestataire mal maîtrisé ouvre une porte que les tests n'ont pas
explorée : la défaillance se propage, et dans un sens. Un domaine en entraîne
d'autres, et il vaut mieux corriger une cause qu'un symptôme. Les outils
statistiques habituels représentent le fait que plusieurs domaines cèdent
ensemble, mais pas lequel a entraîné l'autre, alors que c'est ce sens qui devrait
dicter l'ordre des travaux.

\textbf{Ce qui a été construit.} Le travail représente la propagation par une
\emph{cascade dirigée} : un incident naît sur un domaine, puis se transmet aux
autres selon une matrice d'influence asymétrique, posée à dire d'expert, corroborée par un corpus
de rapports d'incidents officiels et normalisée de sorte que la cascade s'éteigne
toujours. L'état de conformité, domaine par domaine, pilote alors quatre
paramètres de la loi des pertes : combien d'incidents surviennent par an, quelle
proportion échappe à la détection, avec quelle force ils se propagent, et dans
quelle mesure un prestataire commun les fait converger. Le capital est recalculé à
partir de cette loi, sans pénalité ajoutée.

\textbf{Le résultat principal.} Entre un état pleinement conforme et un état
pleinement défaillant, le besoin de capital est multiplié par $3{,}3$. Surtout,
les quatre canaux ne s'additionnent pas : pris séparément ils expliquent les deux
tiers de l'écart, et le tiers restant, environ $35\,\%$, vient de leurs
interactions. Conséquence directe : corriger un domaine rapporte davantage à une
entité défaillante partout qu'à une entité déjà saine par ailleurs, parce que la
correction ferme aussi les interactions que ce domaine portait. Le chiffre à citer
dans un plan d'action est donc l'effet de fermeture, de $1{,}5$ à $3$~fois plus
grand que l'effet isolé.

\textbf{L'apport de méthode, et il est plus important que le chiffre.} La donnée
disponible ne permet pas d'observer le \emph{sens} de la propagation, et ce n'est
pas un manque de volume : la grandeur observable y est mathématiquement
insensible, donc une matrice et son symétrique sont indiscernables pour la donnée
tout en conduisant à des décisions différentes. Un schéma d'incidents standard de
la profession porte les deux champs nécessaires et ne les renseigne ensemble que
dans un incident sur $10\,591$. Le travail ne pose donc pas une valeur
arbitraire : il calcule l'\emph{ensemble} des matrices compatibles avec la donnée
et publie le capital en fourchette. L'ignorance a ainsi un coût connu d'avance et
plafonné, et l'essentiel du surcoût vient de la co-occurrence des défaillances,
qui est observable, non du sens de la propagation. Ce sens manquant déplace la
\emph{priorité} de remédiation plus que le montant, d'où une hiérarchie de
domaines plutôt qu'un chiffre.

\textbf{Les limites, publiées avec le résultat.} Le niveau absolu du capital est
illustratif : il varie fortement selon la base de sinistres retenue, et c'est le
rapport entre les deux états qui est défendu. Un backtest hors échantillon valide
la loi du nombre d'incidents et la forme de la queue des pertes, mais montre que
le niveau des pertes extrêmes dérive plus vite que le modèle ne le suppose, donc
que le capital est sous-estimé. Ce même backtest chiffre enfin sa propre
puissance : détecter une anomalie sur le quantile réglementaire demanderait plus
de mille huit cents années d'observations, donc aucun historique existant ne
permet de valider ce niveau.

\chapter*{Executive summary}
\addcontentsline{toc}{chapter}{Note de synthèse (English)}
\markboth{Executive summary}{}

% HARNAIS-HORS-SECTION: stand-alone executive summary for a non-specialist
% reader. Figures restate, in rounded form, quantities established and verified
% in the body of the memoir, where they carry their script citations.

\textbf{The problem.} Since January 2025, the European DORA regulation has
required banks and insurers to master five areas of digital operational
resilience: technology risk governance, incident management, resilience testing,
third-party provider risk, and cyber threat information sharing. The regulation
states what must be put in place. It does not state what failing to do so costs.
Nor does Solvency~II, the prudential framework that sets an insurer's capital: its
operational risk charge is a percentage of premiums, identical for an exemplary
insurer and for a failing one. An executive weighing a compliance budget therefore
has no figure to work with.

\textbf{The difficulty is not the obvious one.} The five areas are not independent
boxes to tick. Weak governance delays incident detection; an unmanaged provider
opens a door that testing never explored. Failure propagates, and it propagates in
a \emph{direction}: one area drags others down, and fixing a cause is worth more
than fixing a symptom. Standard statistical tools can represent the fact that
several areas fail together, but not which one dragged the others. Yet it is
precisely that direction which should dictate the order of remediation work.

\textbf{What was built.} The work represents propagation as a \emph{directed
cascade}: an incident starts in one area, then transmits to others through an
asymmetric influence matrix posited by expert judgement and corroborated by a corpus of official
post-incident reports. The matrix is normalised so that the cascade always dies out, a property
guaranteed by construction rather than checked afterwards. The firm's compliance
state, area by area, then drives four parameters of the loss distribution: how
many incidents occur per year, what share escapes detection, how strongly they
propagate, and how far a shared provider makes them converge. Capital is recomputed
from that distribution, with no penalty added on top.

\textbf{The headline result.} Between a fully compliant state and a fully failing
one, the capital requirement is multiplied by $3{,}3$. More importantly, the four
channels are not additive: taken separately they account for two thirds of the gap,
and the remaining third, about $35\,\%$, comes from their interactions. One direct
consequence for a remediation plan: fixing one area is worth more to a firm that
fails everywhere than to a firm otherwise sound, because the fix also closes the
interactions that area carried. The figure to quote in an action plan is therefore
not the area's standalone effect but its closure effect, which is $1{,}5$ to $3$
times larger.

\textbf{The methodological contribution, which matters more than the figure.} The
available data does not allow the \emph{direction} of propagation to be observed.
This is not a matter of sample size: the observable quantity is mathematically
insensitive to direction, so that a matrix and its symmetric counterpart are
indistinguishable to the data even though they lead to different decisions. An
industry-standard incident schema carries the two required fields and populates
both in only one incident out of $10\,591$. Faced with this, the work does not
posit an arbitrary value. It computes the \emph{set} of matrices consistent with
the data and reports capital as a range over that set. Ignorance thus has a cost
known in advance and capped, and most of the additional capital comes from the
co-occurrence of failures, which is observable, not from the direction of
propagation, which is not. The missing direction shifts the remediation
\emph{priority} more than the amount, which is the reason for publishing a ranking
of areas rather than a single number.

\textbf{Limitations, published alongside the result.} The absolute capital level is
illustrative: it varies widely with the loss database used, and what is defended is
the ratio between the two compliance states. An out-of-sample backtest validates
the incident count distribution and the shape of the loss tail, but shows that the
level of extreme losses drifts faster than the model assumes, so that capital is
understated. The same backtest finally quantifies its own power: detecting an
anomaly at the regulatory quantile would require more than one thousand eight
hundred years of observations, so no existing history can validate that level.
